\section{The symbolic-rank engine}
\label{sec:symbolic}

Every closed form in \S\ref{sec:second}--\S\ref{sec:grading} is a
rational function of $N$ computed as such, not interpolated through
ranks. This section explains how, because the method is the most
transferable thing in the program and because getting it wrong is the
default.

\subsection{The problem with sweeping ranks}

The natural procedure is to run an exact engine at $N = 3, 4, 5, \ldots$
and fit a rational function. That is what was done first. An assembled
cluster cumulant $\beta_N$ agreed with a candidate closed form
$P_{17}(N^2)\big/\big(N\,R_{20}(N^2)\big)$ at every rank from $4$ to
$70$: sixty-seven exact agreements between exact rationals.

Sixty-seven exact agreements do not prove the identity. Rational
interpolation is unique only once the degrees are bounded, and without
a degree bound the agreements constrain nothing about $N = 71$. The
review of that work made exactly this objection --- no ``for all $N$''
without a degree bound --- and it was correct.

A degree bound could have been argued: Weingarten denominators, energy
denominators, and a count of vertices together bound the degree, and
the sixty-seven agreements would then prove the identity by uniqueness.
That argument was not made. What was done instead makes it unnecessary,
and is better.

\subsection{Running the engine over \texorpdfstring{$\Q(N)$}{Q(N)}}

Inspect where the rank actually enters the exact engine. It enters in
exactly two places that are not rational arithmetic:
\begin{itemize}
  \item the charge predicate, $\text{charge} \equiv 0 \bmod N$;
  \item the flux-difference predicate, $\abs{\mathrm{diff}} \le N$.
\end{itemize}
Everywhere else the rank appears only through rational expressions in
$C_F = (N^2-1)/(2N)$, in $1/(2N)$, and in $N^{\text{closed}}$.

So make the two predicates overridable and promote everything else to
the rational function field:
\begin{itemize}
  \item scalars become elements of $\Q(N)$, carried as a numerator and
    a monic denominator in exact polynomial arithmetic;
  \item the Weingarten function becomes the symbolic inverse of the
    Gram matrix $N^{\text{cycles}}$ on $S_n$;
  \item the single-link spectra become the $\SU(N)$ quadratic Casimirs
    of the irreducible representations that link's flux content can
    carry.
\end{itemize}

A cumulant then emerges as \emph{one rational function of $N$},
computed rather than fitted. The per-rank engine is otherwise untouched
and its own test suite still holds --- which is the check that the
symbolic run is the same calculation and not a different one.

\begin{theorem}[Closed forms over $\Q(N)$]
\label{thm:symbolic}
The fourth-order cluster cumulants, the two-hop weight in both
$C$-parity sectors, and the assembled $\beta_N$ are exact elements of
$\Q(N)$ produced by the symbolic engine. The identity
$\beta_N = P_{17}(N^2)\big/\big(N\,R_{20}(N^2)\big)$ holds as an
identity of rational functions.
\end{theorem}
\gid{ADR:0029}

\noindent The final polynomial identity is promoted to \Tzero{}, and it
is worth saying exactly which identity that is. The corpus \emph{states}
$\beta_N = P_{17}(N^2)/(N R_{20}(N^2))$ and never derives it. What Lean
checks is that the closed forms the symbolic engine returns for the
three cumulants and the adjacent-face cube completion \emph{sum to} the
corpus's coefficient --- a polynomial identity between the engine's
output and the corpus's formula. It is not a Lean derivation of the
cumulants, and the Lean source says so in its own docstring. This
document repeats the qualification because an internal summary of the
same fact dropped it.

Every eigencomponent of every resolvent is verified to be an exact
eigenvector of every single-link $H_0$, symbolically --- a check that
has no per-rank analogue, because at a fixed rank a coincidence of
numbers is indistinguishable from an identity.

\subsection{Why this is the right shape of argument}

The methodological content generalizes beyond this program.

An identity between rational functions can be established in two ways:
verify it at enough points and bound the degrees, or compute in the
function field. The first requires an argument about the size of the
objects, which is a second piece of mathematics and is where the
difficulty hides. The second requires only that the algorithm's
dependence on the parameter be rational --- which is a property one can
\emph{inspect}, predicate by predicate, rather than estimate.

The title of the decision record is the summary: the degree bound is
proved by removing it.

\begin{scopenote}
The symbolic engine is valid at ranks where the generic partition
labelling applies. Exceptional low ranks, where representations
degenerate, are handled separately and are the reason $\SU(2)$ and in
places $\SU(4)$ carry their own entries in the register.

This degree bound must not be confused with the \emph{other} degree
bound in this program's history --- the proposed kinematic mechanism
for the fourth-order tier collapse, which was retracted
(\S\ref{sec:fourth}, \S\ref{sec:register}). That one was withdrawn
because it was wrong. This one was removed because it became
unnecessary.
\end{scopenote}
